\begin{summarybox}
	\textbf{\textcolor{CSummary}{一句话核心}}
	三次函数零点个数由导函数判别式\(\Delta'\)及极值点函数值乘积\(f(x_1)f(x_2)\)的符号共同决定.

	\medskip
	\textbf{\textcolor{CSummary}{使用条件}}
	\begin{itemize}[leftmargin=*, itemsep=0.5em, topsep=0.4em]
		\item \textbf{条件1（函数形式限制）：} 函数必须是三次函数\(f(x)=ax^3+bx^2+cx+d\)，其中\(a\neq0\).
		\item \textbf{条件2（先算判别式）：} 导函数判别式\(\Delta'=4(b^2-3ac)\)必须计算，以判断极值点存在性.
	\end{itemize}

	\medskip
	\textbf{\textcolor{CSummary}{关键提醒}}
	\begin{itemize}[leftmargin=*, itemsep=0.5em, topsep=0.4em]
		\item \textbf{易错点（乘积符号记忆）：} 牢记口诀：负三零，正一零，零二重（重根）。
		\item \textbf{检查项（前提条件）：} 使用乘积符号判定前，务必确认\(\Delta'>0\)；否则函数单调，只有一个零点.
	\end{itemize}

\end{summarybox}
